% !TEX root = ../main.tex

\subsection{Dataset}

We use the Impurities in 2D Materials Database
(IMP2D)~\cite{imp2d2023}, filtering for DFT-converged configurations with
$|\ef| < 20$\eV{}, yielding 10\,641 samples spanning 44 host monolayers and
65 dopant elements.
The dataset covers two defect types: adsorbates (68\%) and interstitials
(32\%).
We adopt a deterministic ordered split of 8\,512 / 1\,064 / 1\,065 for
training / validation / testing.

For multi-source training, we additionally incorporate JARVIS-2D (2D
formation energies), JARVIS-3D (3D formation energies), and a curated DFT-3D
set, totaling ${\sim}27$\,k samples across all sources.

\subsection{Baselines}

\begin{itemize}
  \item \textbf{Naive baseline}: global mean of training-set $\ef$ values,
    providing a lower bound on model utility.
  \item \textbf{LightGBM}: tree-based baseline with 9-dimensional atomic
    features (same as \dart{} node features), providing a non-neural reference.
  \item \textbf{El~Alouani~\textit{et al.}}~\cite{elalouani2024}:
    LightGBM with Jacobi--Legendre structural descriptors on the same
    IMP2D dataset (reported result).
  \item \textbf{SchNet}~\cite{schutt2018}: continuous-filter GNN
    (0.46\,M parameters), representing the CGCNN-class of models without
    angular information.
  \item \textbf{ALIGNN}~\cite{alignn2021}: line-graph GNN with explicit
    bond-angle encoding (4.03\,M parameters), the current standard for
    crystal property prediction.
  \item \textbf{CrystalTransformer}~\cite{crystalformer2024}: the vanilla
    local+global transformer with distance-biased attention
    (0.75\,M parameters), trained with MSE loss and plateau learning rate
    scheduling.
    \dart{} shares this backbone architecture but adds the defect-site
    embedding, switches to an optimized training recipe
    (\S\ref{sec:training}), and incorporates multi-source fusion---making
    the comparison between the two a direct measure of these contributions.
\end{itemize}

All GNN baselines are trained with identical data splits and evaluation
protocol for fair comparison.

\subsection{Training details}

\dart{} is trained with AdamW ($\beta_1 = 0.9$, $\beta_2 = 0.999$,
weight decay $10^{-4}$), batch size 64, and gradient clipping at 5.0.
Host-balanced sampling ensures each of the 44 hosts contributes equally per
epoch (200 samples per host with replacement).
All experiments use a single NVIDIA L40S GPU; typical training time is
${\sim}22$ minutes for 150 epochs.

For the OOD evaluation (Tier~1), we retrain \dart{} from scratch for each of
the 7 LOHO folds, using the standard recipe without multi-source data
(single-source IMP2D only, seed 42).

\subsection{Evaluation metrics}

We report mean absolute error
$\mae = \frac{1}{N}\sum_i |{\hat{E}}_{\mathrm{f},i} - E_{\mathrm{f},i}|$
and root-mean-square error (RMSE) in eV.
For OOD experiments, we additionally compare against the naive baseline to
quantify the improvement ratio.
Statistical significance of ablation comparisons is assessed via paired
$t$-tests across multiple seeds (3--5 seeds per configuration).
